\documentclass[10pt,a4paper]{amsart}
\usepackage{amsmath,amssymb,amsthm,graphicx}
% For commutative diagrams we use Xy-pic
\usepackage[all]{xy}

% We use 2cell for 2-commutative diagrams.
\xyoption{2cell}
% Theorem environments.
%
\theoremstyle{plain}
\newtheorem{theorem}[subsection]{Theorem}
\newtheorem{proposition}[subsection]{Proposition}
\newtheorem{lemma}[subsection]{Lemma}

\theoremstyle{definition}
\newtheorem{definition}[subsection]{Definition}
\newtheorem{example}[subsection]{Example}
\newtheorem{exercise}[subsection]{Exercise}
\newtheorem{situation}[subsection]{Situation}

\theoremstyle{remark}
\newtheorem{remark}[subsection]{Remark}
\newtheorem{remarks}[subsection]{Remarks}

\numberwithin{equation}{subsection}

% Macros
%
\def\lim{\mathop{\mathrm{lim}}\nolimits}
\def\colim{\mathop{\mathrm{colim}}\nolimits}
\def\Spec{\mathop{\mathrm{Spec}}}
\def\Hom{\mathop{\mathrm{Hom}}\nolimits}
\def\Ext{\mathop{\mathrm{Ext}}\nolimits}
\def\SheafHom{\mathop{\mathcal{H}\!\mathit{om}}\nolimits}
\def\SheafExt{\mathop{\mathcal{E}\!\mathit{xt}}\nolimits}
\def\Sch{\mathit{Sch}}
\def\Mor{\mathop{\mathrm{Mor}}\nolimits}
\def\Ob{\mathop{\mathrm{Ob}}\nolimits}
\def\Sh{\mathop{\mathit{Sh}}\nolimits}
\def\NL{\mathop{N\!L}\nolimits}
\def\CH{\mathop{\mathrm{CH}}\nolimits}
\def\proetale{{pro\text{-}\acute{e}tale}}
\def\etale{{\acute{e}tale}}
\def\QCoh{\mathit{QCoh}}
\def\Ker{\mathop{\mathrm{Ker}}}
\def\Im{\mathop{\mathrm{Im}}}
\def\Coker{\mathop{\mathrm{Coker}}}
\def\Coim{\mathop{\mathrm{Coim}}}

% Boxtimes
%
\DeclareMathSymbol{\boxtimes}{\mathbin}{AMSa}{"02}

%
% Macros for moduli stacks/spaces
%
\def\QCohstack{\mathcal{QC}\!\mathit{oh}}
\def\Cohstack{\mathcal{C}\!\mathit{oh}}
\def\Spacesstack{\mathcal{S}\!\mathit{paces}}
\def\Quotfunctor{\mathrm{Quot}}
\def\Hilbfunctor{\mathrm{Hilb}}
\def\Curvesstack{\mathcal{C}\!\mathit{urves}}
\def\Polarizedstack{\mathcal{P}\!\mathit{olarized}}
\def\Complexesstack{\mathcal{C}\!\mathit{omplexes}}
% \Pic is the operator that assigns to X its picard group, usage \Pic(X)
% \Picardstack_{X/B} denotes the Picard stack of X over B
% \Picardfunctor_{X/B} denotes the Picard functor of X over B
\def\Pic{\mathop{\mathrm{Pic}}\nolimits}
\def\Picardstack{\mathcal{P}\!\mathit{ic}}
\def\Picardfunctor{\mathrm{Pic}}
\def\Deformationcategory{\mathcal{D}\!\mathit{ef}}
\usepackage{cleveref}
\crefname{originalreference}{original reference}{original references}
\Crefname{originalreference}{Original reference}{Original references}
\makeatletter
\newlabel{lemma-base-change-Rf-star-colim}{{0EZT}{source}{}{sourcepointer-1}{}}
\newlabel{lemma-base-change-Rf-star-colim@cref}{{[originalreference][0][]0EZT}{[1][0][]source}{}{}{}}
\makeatother
\crefname{originalreference}{original reference}{original references}
\Crefname{originalreference}{Original reference}{Original references}
\makeatletter
\newlabel{lemma-base-change-f-star-general}{{0EZX}{source}{}{sourcepointer-2}{}}
\newlabel{lemma-base-change-f-star-general@cref}{{[originalreference][0][]0EZX}{[1][0][]source}{}{}{}}
\makeatother
\crefname{originallemma}{original lemma}{original lemmas}
\Crefname{originallemma}{Original lemma}{Original lemmas}
\makeatletter
\newlabel{varieties-lemma-separably-closed-field-connected-components}{{0363}{source}{}{sourcepointer-3}{}}
\newlabel{varieties-lemma-separably-closed-field-connected-components@cref}{{[originallemma][0][]0363}{[1][0][]source}{}{}{}}
\makeatother
\title{Base change for a pushforward of a sheaf of sets over the spectrum of a separably closed field}
\author{The Stacks Project Authors}
\date{stacks-project:etale-cohomology.tex:local-snapshot-20261011; GNU FDL 1.2 or later}
\begin{document}
\maketitle
\noindent\textit{Editorial note.} All delivered mathematical text is the author's own, copied from the pinned Stacks Project chapter file etale-cohomology.tex; no statement or proof step was rewritten, reordered or removed.\par\medskip
\noindent\textit{Editorial note.} The counted proof is the author's complete argument: the proof is written out in full: the localisation and the cofiltered limit presentation, the reduction to finite type, and the application of the general base change lemma with the geometric connectedness input.\par\medskip
\noindent\textit{Editorial note.} The label lemma-base-change-f-star-field names the counted statement, and the proof environment immediately following it in the source is delivered verbatim as its proof.\par\medskip
\noindent\textit{Editorial note.} The counted statement is carried by the author's own lemma environment in the chapter source.\par\medskip
\noindent\textit{Editorial note.} The statement's cartesian diagram is the author's and is delivered verbatim; the three cited inputs are kept as external pointers with their canonical Stacks tags.\par\medskip
\noindent\textit{Allowed standard prerequisites (editorial).}\begin{itemize}
\item Base change for a filtered colimit of schemes. Same chapter, the lemma comparing Rg\_* of a colimit with the colimit of the comparisons, used to reduce to finite type over S. Supplies the passage from an affine X presented as a cofiltered limit to a finite type model.
\item Base change in the presence of geometrically connected fibres. Same chapter, the general base change lemma for a flat morphism whose strict henselisations have geometrically connected fibres. Supplies the conclusion once the fibres are known to be geometrically connected.
\item Structure of schemes of finite type over a separably closed field. Varieties, the lemma stating that such a scheme is a finite disjoint union of connected and geometrically connected schemes. Supplies the geometric connectedness of the fibres of the finite type model.
\end{itemize}\medskip
\ifdefined\hypertarget\hypertarget{sourcepointer-1}{}\fi
\noindent\textit{Source reference.} Original reference 0EZT: sources/stacks-etale-cohomology/pinned-original-source.tex, source lines 16410--16410. Prerequisite pointer invoked at proof line(s) 16701; it names a standard fact of the same chapter that the counted argument uses, and every step of the counted proof built on it is retained verbatim in the delivered body.\par\medskip
\ifdefined\hypertarget\hypertarget{sourcepointer-2}{}\fi
\noindent\textit{Source reference.} Original reference 0EZX: sources/stacks-etale-cohomology/pinned-original-source.tex, source lines 16610--16610. Prerequisite pointer invoked at proof line(s) 16703; it names a standard fact of the same chapter that the counted argument uses, and every step of the counted proof built on it is retained verbatim in the delivered body.\par\medskip
\ifdefined\hypertarget\hypertarget{sourcepointer-3}{}\fi
\noindent\textit{Source reference.} Original lemma 0363: sources/stacks-chow/tags.txt, source lines 4150--4150. Prerequisite pointer invoked at proof line(s) 16708; it names a standard fact of a sibling chapter that the counted argument uses, and every step of the counted proof built on it is retained verbatim in the delivered body.\par\medskip
\begin{lemma}
\label{lemma-base-change-f-star-field}
Consider the cartesian diagrams of schemes
$$
\xymatrix{
X \ar[d]_f & Y \ar[l]^h \ar[d]^e \\
S & T \ar[l]_g
}
$$
Assume that $S$ is the spectrum of a separably closed field.
Then $f^{-1}g_*\mathcal{F} = h_*e^{-1}\mathcal{F}$
for any sheaf $\mathcal{F}$ on $T_\etale$.
\end{lemma}

\begin{proof}
We may work locally on $X$. Hence we may assume $X$ is affine.
Then we can write $X$ as a cofiltered limit of affine schemes of
finite type over $S$. By Lemma \ref{lemma-base-change-Rf-star-colim}
we may assume that $X$ is of finite type over $S$.
Then Lemma \ref{lemma-base-change-f-star-general}
applies because any scheme of finite
type over a separably closed field is a finite disjoint
union of connected and geometrically connected schemes
(see Varieties, Lemma
\ref{varieties-lemma-separably-closed-field-connected-components}).
\end{proof}

\end{document}
